\documentclass[11pt,a4paper]{article}
\usepackage[utf8]{inputenc}
\usepackage[T1]{fontenc}
\usepackage[american,polish]{babel}
\usepackage{lmodern}
\usepackage[margin=27mm]{geometry}
\usepackage{microtype}
\usepackage{amsmath}
\usepackage{csquotes}
\usepackage[backend=biber,style=apa]{biblatex}
\DeclareLanguageMapping{polish}{american-apa}
\addbibresource{references.bib}
\usepackage[hidelinks]{hyperref}
\title{Odbiór rzeczywistych zachowań osób neuroróżnorodnych po przeniesieniu na syntetyczne awatary}
\author{Mateusz Pomianek}
\date{Szkic naukowy, wersja 0.3.0, 27 listopada 2026}
\begin{document}
\maketitle
\begin{abstract}
Czy rzeczywiste wzorce komunikacji osób o różnych profilach, odtwarzane na identycznych awatarach, wywołują odmienne pierwsze wrażenie zależnie od wyglądu postaci? Przedstawiono prospektywny projekt badania oparty na wspólnej platformie ST-08, wykorzystujący archiwalne nagrania zachowań odtwarzane na syntetycznych postaciach o kontrolowanym wyglądzie. Pierwszorzędowy wskaźnik: Zależność ocen niechęci i niesamowitości od indywidualnych parametrów zachowania oraz realizmu postaci z uwzględnieniem kontekstu. Badanie nie dostarczyło dotąd danych ani wyników empirycznych.
\end{abstract}
\section{Wprowadzenie}
Zależność między wyglądem jednostki a widocznym przebiegiem jej zachowania jest zasadniczym elementem społecznej percepcji. W kontaktach z robotami społecznymi i realistycznymi awatarami nieoczekiwana mimika, gest, kontakt wzrokowy i dynamika wypowiedzi mogą wpływać na reakcje odbiorcy. Nie oznacza to jednak, że każda atypowa ekspresja musi być oceniana negatywnie, ani że jest uprawnione utożsamianie ekspresji z intensywnością doświadczanych emocji.

Projekt stanowiska ST-08 umożliwi oddzielenie zachowania rzeczywistego człowieka, wyodrębnionego z istniejących nagrań interakcji, od syntetycznej formy jego reprezentacji. Pozwala to porównywać te same zdarzenia w kontrolowanym wyglądzie awatara. Nie można jednak przyjąć, że retargeting automatycznie zachowuje wszystkie własności zachowania. Jego dokładność i sposób deformowania sygnałów muszą być sprawdzone niezależnie od właściwego eksperymentu.
\section{Podstawy teoretyczne}
Oceny formułowane na podstawie krótkich fragmentów zachowania stanowią reakcję obserwatora, nie bezpośredni pomiar kompetencji osoby źródłowej. \textcite{sasson2017} wykazali zależność pierwszych wrażeń od dostępności audiowizualnych informacji o zachowaniu w badaniach osób autystycznych. Wynik ten uzasadnia kontrolę kanałów prezentacji, ale nie dowodzi, że latencja odpowiedzi odpowiada za obserwowane różnice.

Koncepcja podwójnej empatii ujmuje trudności komunikacyjne jako potencjalnie wzajemne \parencite{milton2012}. Wyniki badań nad jakością kontaktu w parach o zgodnych i różnych neurotypach wskazują na znaczenie relacyjnego kontekstu \parencite{crompton2020}. Oceny pojedynczego klipu nie są jednak pomiarem wzajemnej empatii ani jakości rzeczywistej interakcji. Uogólnienie badań autyzmu na wszystkie profile neuroróżnorodności wymaga odrębnych danych.

Przegląd \textcite{katsyri2015} dostarcza przesłanek na rzecz niedopasowania percepcyjnego, lecz nie uzasadnia uniwersalnej zależności między humanizacją awatara i dyskomfortem. W projekcie P04 realizm jest eksperymentalnym moderatorem naturalnie występujących cech komunikacji, nie narzędziem klasyfikowania osób. Niesamowitość i gotowość do kontaktu są odrębnymi konstruktami. Alternatywne wyjaśnienia obejmują artefakty animacji, atrakcyjność wyglądu, znajomość awatarów oraz właściwości sytuacji komunikacyjnej.

\section{Cel i hipotezy}
Czy rzeczywiste wzorce komunikacji osób o różnych profilach, odtwarzane na identycznych awatarach, wywołują odmienne pierwsze wrażenie zależnie od wyglądu postaci?
\paragraph{H1.} Czasowe i przestrzenne właściwości naturalnie zarejestrowanej ekspresji komunikacyjnej są związane z ocenami pierwszego wrażenia przy kontrolowanej tożsamości awatara.
\paragraph{H2.} Realizm wizualny moderuje związek między indywidualnym wzorcem zachowania a oceną niechęci i niesamowitości.
Przed badaniem potwierdzającym konieczne jest ustalenie estymandu, sposobu kodowania efektu, kryteriów jakości, mocy oraz procedury radzenia sobie z wielokrotnymi porównaniami.
\section{Materiał i metody}
Archiwalne interakcje z legalnym prawem do retargetingu, pełnym kontekstem i niezależnie udokumentowanymi informacjami o osobie źródłowej. Rozpoznania kliniczne nie są inferowane z wideo.

Wielu źródłowych uczestników i wiele kontekstów. Przeniesienie przebiegów mimiki, gestów i dynamiki spojrzenia na tę samą rodzinę awatarów; oceniający nie otrzymują etykiet klinicznych.

Pierwszorzędowy wskaźnik: Zależność ocen niechęci i niesamowitości od indywidualnych parametrów zachowania oraz realizmu postaci z uwzględnieniem kontekstu. Jednostką analizy jest ocena konkretnej sekwencji przez konkretnego obserwatora; należy uwzględnić zależność prób od źródłowej osoby, scenariusza i postaci. Zbieranie odpowiedzi behawioralnych nie wymaga EEG ani okulografii. Szacowana orientacja spojrzenia w archiwalnym filmie nie jest laboratoryjnym pomiarem eye trackingu.
\section{Plan analizy}
Hierarchiczny model obserwatora, źródłowej osoby i scenariusza. Ocena rzetelności ekstrakcji, artefaktów i zgodności warunków; brak uogólniania na diagnozę bez odpowiedniej próby.

Zamrożenie estymandu i planu porównań nastąpi po pilotażu technicznym i przed dostępem do danych konfirmacyjnych. Należy zachować oddzielnie analizy eksploracyjne, raportować przedziały niepewności, wielkości efektów oraz reguły wyłączenia bodźców z błędną animacją. Trafność konstruktu wymaga niezależnej kalibracji realizmu wyglądu i percepcyjnej naturalności zachowania.
\section{Projekt konfirmacyjny}
Proponuje się trzy poziomy realizmu wizualnego, kalibrowane niezależnie, oraz powtarzanie tych samych sekwencji źródłowych w różnych reprezentacjach awatara. Obserwator otrzymuje tylko jedną wersję danej sekwencji, a przydział wersji jest randomizowany i równoważony między obserwatorami. Każda osoba źródłowa dostarcza wiele sekwencji w porównywalnych kontekstach. Tożsamości awatarów są skrzyżowane z osobami źródłowymi i realizmem. Jednostką niezależnej informacji o cechach zachowania pozostaje osoba źródłowa, a nie liczba ocen.

Jako roboczy predyktor główny proponuje się medianę latencji odpowiedzi gestem lub ruchem głowy po jednoznacznie oznaczonym zdarzeniu komunikacyjnym. Okno pomiarowe, reguły kodowania i postępowanie z brakiem odpowiedzi zostaną ustalone w pilotażu. Predyktor nie jest wskaźnikiem diagnozy ani deficytu. Sekwencje bez obserwowalnego zdarzenia nie dostarczają tego pomiaru; ich udział należy raportować. Rzetelność kodowania i zgodność pomiaru przed retargetingiem oraz po nim wymagają niezależnej oceny.

Pierwszorzędowym punktem końcowym proponuje się pięciokategorialną ocenę niesamowitości; gotowość do interakcji pozostaje odrębnym wynikiem drugorzędowym. Dla standaryzowanego predyktora $x$ oraz kontrastu wysokiego i niskiego realizmu $r$ model kumulacyjny ma postać
\[
\operatorname{logit}P(Y\leq c)=\theta_c-(\beta_x x+\beta_r r+\beta_{xr}xr+\boldsymbol{\gamma}^{T}\mathbf{z}+u_o+u_s+u_k+u_a).
\]
Wektor $\mathbf{z}$ obejmuje kontrast pośredniego realizmu, jego interakcję z predyktorem oraz kontekst. Efekty losowe dotyczą obserwatora, osoby źródłowej, klipu zagnieżdżonego w osobie i awatara. Należy rozważyć losowe nachylenia realizmu dla obserwatora i osoby źródłowej. H2 odpowiada współczynnikowi $\beta_{xr}$; H1 pozostaje hipotezą drugorzędową. Randomizacja realizmu umożliwia interpretację jego efektu przyczynowego w granicach projektu, ale związek z naturalnie występującą cechą zachowania pozostaje obserwacyjny.

Symulacja Monte Carlo powinna niezależnie zmieniać liczbę osób źródłowych, obserwatorów, klipów, błąd pomiaru predyktora i wariancję efektów losowych. Analiza wrażliwości obejmie selekcję klipów i brak pomiaru latencji. Nie ustalono liczebności docelowej.

Przed rekrutacją konieczne są potwierdzenie praw do retargetingu i przetwarzania danych, kwalifikacja etyczna oraz ocena ryzyka stygmatyzacji z udziałem konsultantów neuroróżnorodnych. Nie wolno publikować identyfikowalnych nagrań ani wykorzystywać ocen do klasyfikacji klinicznej osób źródłowych.
\section{Walidacja bodźców}
Procedura obejmuje niezależne kodowanie latencji, ocenę zgodności czasowej przed retargetingiem i po nim, kontrolę artefaktów oraz niezależny pilotaż realizmu i wierności ruchu. Dobór klipów nie może zależeć od ich ocen niesamowitości. Kryteria akceptacji, marginesy równoważności i reguły wyłączeń wymagają zamrożenia przed badaniem głównym. Szczegółowy protokół walidacji pozostaje dokumentem prospektywnym.
\section{Wyniki}
Nie zebrano jeszcze danych empirycznych i nie wykonano analiz potwierdzających. W niniejszym opracowaniu nie przypisuje się wartości statystycznych, wyników skal, skuteczności diagnostycznej ani dowodów występowania doliny niesamowitości.
\section{Dyskusja i ograniczenia}
Naturalna neuroróżnorodność nie jest skalą mechaniczności ruchu. Badamy zachowania osób źródłowych, a nie profil neuroróżnorodności obserwatorów. Prawa do korzystania z oryginalnych nagrań i rozpowszechniania pochodnych animacji nie wynikają automatycznie z publicznej dostępności bazy. Rekonstrukcja mimiki i ruchu z pojedynczej kamery może zawierać systematyczne błędy. Dla badań z udziałem nowych obserwatorów konieczne są właściwa kwalifikacja etyczna, ochrona danych i jasno określona procedura zgody.

Wyniki negatywne oraz niejednoznaczne mają znaczenie naukowe, jeśli kontrolowano jakość bodźców, precyzję analizy i heterogeniczność kontekstu. Szczególnie w przypadku różnorodności komunikacji nie można wywodzić wewnętrznych stanów psychicznych z nietypowej ekspresji widocznej na twarzy ani sprowadzać zróżnicowania do jednolitej cechy klinicznej.
\section{Wnioski}
Zaproponowano odrębny problem naukowy i sposób jego przyszłego falsyfikowania. Uprawnienia do nagrań, walidacja rekonstrukcji zachowania, prerejestracja i badania obserwatorów pozostają do wykonania. Powstanie szkicu LaTeX nie stanowi dowodu gotowości publikacyjnej.
\printbibliography
\end{document}
